\subsection{Coupled cyclic rotations and fusion integrality}
\label{sec:m5-cyclic-rotation}
The rotation constraints below couple indicators of a simple object with
indicators of its tensor powers.  They are uniform necessary conditions;
the subsequent finite-domain argument must separately justify its source
dimensions, phase inventory and exhaustive coverage.

\begin{proposition}[Coupled cyclic multiplicities]
\label{prop:m5-cyclic-multiplicities}
Let $X$ be an object of a pivotal fusion category over $\mathbb C$, and
put $h_n(X)=\dim\operatorname{Hom}(\mathbf1,X^{\otimes n})$.
There exist nonnegative integers $u_0,\ldots,u_3$ and
$v_0,\ldots,v_7$ satisfying
\begin{align*}
 \sum_{j=0}^3u_j&=h_4(X),&
 \sum_{j=0}^3i^ju_j&=\nu_4(X),&
 \sum_{j=0}^3(-1)^ju_j&=\nu_2(X^{\otimes2}),\\
 \sum_{j=0}^7v_j&=h_8(X),&
 \sum_{j=0}^7\zeta_8^jv_j&=\nu_8(X),\\
 \sum_{j=0}^7i^jv_j&=\nu_4(X^{\otimes2}),&
 \sum_{j=0}^7(-1)^jv_j&=\nu_2(X^{\otimes4}).
\end{align*}
For a simple $X_i$, the same indicators occurring in these equations
satisfy
\[
 \nu_4(X_i^{\otimes2})=\sum_kN_{ii}^k\nu_4(X_k),\qquad
 \nu_2(X_i^{\otimes r})=\sum_k[X_i^{\otimes r}:X_k]\nu_2(X_k)
 \quad(r=2,4).
\]
\end{proposition}
\begin{proof}
Let $E_n$ be pivotal rotation on
$\operatorname{Hom}(\mathbf1,X^{\otimes n})$.
The block-composition identity
\cite[equation~(5.4)]{NgSchauenburgPivotal} identifies $E_n^r$,
for $r\mid n$, with rotation of the $r$-fold blocks.
Moreover $E_n^n=1$ by \cite[Theorem~5.1]{NgSchauenburgPivotal}.
Thus $u_j$ and $v_j$ can be taken as the eigenspace dimensions for
eigenvalues $i^j$ and $\zeta_8^j$. Taking traces of the indicated powers
gives the first identities. Additivity of indicators gives the last two.
\end{proof}

For explicit use, write $\nu_4(X)=a+ib$, $m=h_4(X)$ and
$r=\nu_2(X^{\otimes2})$.  The four required multiplicities are
\[
 \frac14(m+r+2a,\ m-r+2b,\ m+r-2a,\ m-r-2b).
\]
For the eighth rotation put $m=h_8(X)$, $q=\nu_2(X^{\otimes4})$,
$\nu_4(X^{\otimes2})=s+it$, and write
$\nu_8(X)=\sum_{j=0}^3\eta_j\zeta_8^j$.  Set
\[
 A_0=\frac{m+q+2s}{4},\quad A_2=\frac{m+q-2s}{4},\quad
 A_1=\frac{m-q+2t}{4},\quad A_3=\frac{m-q-2t}{4}.
\]
Then $v_j=(A_j+\eta_j)/2$ and $v_{j+4}=(A_j-\eta_j)/2$.
Requiring these eight values to be nonnegative integers gives the exact
cyclic constraint. In the basis $1,\sqrt2,i,i\sqrt2$, coefficients
of $\sqrt2$ and $i\sqrt2$ can be half-integral; imposing integral
coordinates in that basis would add an invalid restriction.

\begin{proposition}[Extraction from center restrictions]
\label{prop:m5-cyclic-center-trace}
Assume now that the category is spherical. Let $F$ be the forgetful
functor from its center, let $D$ be its global dimension, and write
$d_i=\dim(X_i)$, $v_{Yi}=[F(Y):X_i]$ and
$z_Y=\sum_i v_{Yi}d_i$ for a central simple $Y$.
Choose the center twist convention of
\cite[Theorem~4.1]{NgSchauenburgExponent}.  Then
\[
 D\nu_n(X_i)=\sum_Yv_{Yi}z_Y\theta_Y^n.
\]
If $K$ contains the source dimensions and $\zeta_8$, normalized relative
trace to $K$ therefore gives the complete fourth and eighth indicators:
\[
 D\nu_n(X_i)=\sum_Yv_{Yi}z_Y\,t_K(\theta_Y^n),\qquad n=4,8.
\]
\end{proposition}
\begin{proof}
The first identity is the center-indicator formula. Its relative trace
commutes with the coefficients $v_{Yi}z_Y\in K$.
Proposition~\ref{prop:m5-cyclic-multiplicities} places $\nu_n(X_i)$
in $\mathbb Q(\zeta_n)\subset K$ for $n=4,8$, so the left side is
fixed by that trace. The opposite center convention replaces every twist
by its inverse, as in \cite[Remark~4.3]{NgSchauenburgExponent}; changing
conventions uniformly permutes the cyclic eigenvalue labels.
\end{proof}

The restriction Gram identity, fusion tensor and complete relative-phase
inventory now give rational linear equations in nonnegative integral
restriction and cyclic multiplicities. Clearing denominators produces an
integer system. A modular row functional, an exact rational separator,
or an exhaustive finite branch certificate may prove its infeasibility.
This conclusion requires verification of the whole finite domain and every
dimension/indicator branch. A general complex exclusion additionally uses
the independently established reduction from arbitrary categorifications
to the spherical branches under consideration.

\begin{proposition}[Integral fusion-norm inequality]
\label{prop:m5-integral-fusion-norm}
Let $A,B$ be central simple objects and put $n_Y=[A\otimes B:Y]$,
$a=[F(A)\otimes F(B)]$ and $q=\dim\operatorname{End}(A\otimes B)$.
Partition the central simples into a specified low set $\mathcal L$
and its complement $\mathcal H$. For a rational source-coordinate
vector $x$ define
\[
 g^*(t)=\max_{k\in\mathbb Z_{\geq0}}(tk-k^2).
\]
For arbitrary integers $b_G\geq0$, $G\in\mathcal L$, one has
\begin{align*}
 q\ \geq\ x\cdot a
 &+\sum_{G\in\mathcal L}(2b_G+1-x\cdot v_G)n_G
   -\sum_{G\in\mathcal L}b_G(b_G+1)\\
 &-\sum_{Y\in\mathcal H}g^*(x\cdot v_Y).
\end{align*}
\end{proposition}
\begin{proof}
Semisimplicity and restriction give $q=\sum_Yn_Y^2$ and
$a=\sum_Yn_Yv_Y$. For $Y\in\mathcal H$, the definition of $g^*$ gives
$n_Y^2\geq(x\cdot v_Y)n_Y-g^*(x\cdot v_Y)$.
For $G\in\mathcal L$, the product of the consecutive integers
$n_G-b_G$ and $n_G-b_G-1$ is nonnegative, so
$n_G^2\geq(2b_G+1)n_G-b_G(b_G+1)$. Summing proves the inequality.
The maximum defining $g^*$ is computed exactly by checking
$\max(0,\lfloor t/2\rfloor)$ and the next integer.
\end{proof}

For $A,B$ in a fixed labelled low frame, $a$ and the low restriction
rows are fixed. The low coefficients $n_G$ and $q$ are respectively triple and
quartic Verlinde moments; the final sum is linear in high-row
multiplicities. Thus fixed rational choices of $x$ and integral choices
of $b_G$ produce valid linear inequalities for the complete necessary
system. A failed trial witness may suggest those choices, but the proof
of the inequality does not depend on that witness or on an optimizer.

\begin{proposition}[All-order power-of-two phase inventory]
\label{prop:m5-power-phase-inventory}
Put $K=\mathbb Q(\sqrt5,\zeta_8)$ and let $k\geq1$.
For roots of unity $\theta$ whose orders have no prime factors other than
$2$ and $5$, the vectors
\[
 \bigl(t_K(\theta),t_K(\theta^2),\ldots,
             t_K(\theta^{2^k})\bigr)
\]
take exactly $25+12k$ values. All nonzero values are realized at orders
dividing $40\cdot2^k$. The zero vector corresponds precisely to exact
orders $2^a5^b$ with $a\geq k+4$ or $b\geq2$.
\end{proposition}
\begin{proof}
The two-primary trace vanishes when $a>j+3$ at the $2^j$-th power;
otherwise the powered root lies in $\mathbb Q(\zeta_8)$ and is unchanged
by the trace. For $b\geq2$ every indicated five-primary trace vanishes.
For $b=0$ the five-primary factor is $1$, whereas for $b=1$ its successive
traces alternate between
$(-1+\delta\sqrt5)/4$ and $(-1-\delta\sqrt5)/4$, with $\delta=\pm1$.
The cyclotomic Chinese remainder decomposition makes the two factors
independent. If the first nonzero coordinate has index $h=0$, there are
eight eighth-root phases and three five-primary choices, giving $24$
vectors. If $1\leq h\leq k$, the first nonzero eighth-root phase is
primitive, giving four choices and hence $12$ vectors for each $h$.
These vectors are distinct, and adjoining zero proves the count.
\end{proof}

For $k=3$ this gives $61$ records for moments $1,2,4,8$: $60$ visible
records represented at orders dividing $320$, and one genuine zero record
covering arbitrarily large orders. Thus $320$ is a modulus for the visible
trace values, rather than an asserted bound on all twists.
The prime restriction used in this application comes from the separately
proved categorical Cauchy argument, not from cyclic rotation alone.

Packet divisibility uses a field containing both source-character fields,
$L=\mathbb Q(\sqrt5,\sqrt{-2})$, which is smaller than the trace field $K$.
It equals the source-character field for $R7\_5\_5039$; that of
$R7\_5\_3895$ is $\mathbb Q(\sqrt5)$. Thus the $L$-fixing subgroup
is a valid, possibly smaller subgroup for either ring.
The square action of automorphisms fixing $L$ has a primitive-root orbit
of size $2^{\max(a-3,0)}5^{\max(b-1,0)}$ at order $2^a5^b$:
its two-primary image is the principal subgroup congruent to $1$ modulo
$8$, and its five-primary image is congruent to $1$ modulo $5$.
Consequently the zero record has per-row multiplicity in
$16\mathbb Z_{\geq0}+5\mathbb Z_{\geq0}$ for $k=3$.
As in the earlier transport argument, restriction rows are fixed only
after the actual positive source-field Galois action and its signs have
been justified. No unsigned raw transport assertion is used here.
